\subsection{Selection of Motion Primitive}\label{sec:methods-selection}

As motivated in \Cref{sec:introduction}, robot capability, object characteristics, and environment configuration inform the need for a diverse repertoire of prehensile and nonprehensile motion primitives.
In this section, we provide a functional definition of our chosen primitive of \textsl{poking}.

\textsl{Poking} is a fundamental nonprehensile motion primitive that is composed of two sequential phases: i) \textsl{impact}: a robot end-effector applies an instantaneous force to an object at rest, setting it in planar translational and rotational motion; and, ii) \textsl{free-sliding}: the object slows down and comes to a halt due to Coulomb friction between the object and its support surface \cite{pasricha2022pokerrt}. This primitive is not limited by robot reach and leads to faster manipulation of objects. However, poking is limited by robot dynamics and nonlinearities due to friction between the object and its support surface \ap{cite}.

